\olchapter{fol}{ntd}{Natural Deduction}

\begin{editorial}
This chapter presents a natural deduction system in the style of
Gentzen/Prawitz.
  
To include or exclude material relevant to natural deduction as a
proof system, use the ``prfND'' tag.
\end{editorial}



\olfileid{fol}{ntd}{rul}

\olsection{Rules and \usetoken{P}{derivation}}

\begin{explain}
Natural deduction systems are meant to closely parallel the informal
reasoning used in mathematical proof (hence it is somewhat
``natural''). Natural deduction proofs begin with assumptions.
Inference rules are then applied. Assumptions are ``!!{discharged}''
by the \Intro{\lnot}, \Intro{\lif}, \Elim{\lor} and
  \Elim{\lexists} inference rules, and the label of
the !!{discharged} assumption is placed beside the inference for
clarity.
\end{explain}

\begin{defn}[Assumption]
An \emph{assumption} is any !!{sentence}
in the topmost position of any branch.
\end{defn}

!!^{derivation}s in natural deduction are certain trees of
!!{sentence}s, where the topmost !!{sentence}s are assumptions, and if
!!a{sentence} stands below one, two, or three other sequents, it must
follow correctly by a rule of inference. The !!{sentence}s at the top
of the inference are called the \emph{premises} and the !!{sentence}
below the \emph{conclusion} of the inference.  The rules come in
pairs, an introduction and an elimination rule for each
!!{operator}. They introduce !!a{operator} in the conclusion or
remove !!a{operator} from a premise of the rule.  Some of the rules
allow an assumption of a certain type to be \emph{!!{discharged}}. To
indicate which assumption is !!{discharged} by which inference, we
also assign labels to both the assumption and the inference.  This is
indicated by writing the assumption as ``$\Discharge{!A}{n}$.''



It is customary to consider rules for all the !!{operator}s $\land$, $\lor$, $\lif$, $\lnot$, and $\lfalse$, even if some of those are defined.







\olfileid{fol}{ntd}{prl}

\olsection{Propositional Rules}


\subsection{Rules for $\land$}

\begin{defish}
\AxiomC{$!A$}
\AxiomC{$!B$}
\RightLabel{\Intro{\land}}
\BinaryInfC{$!A \land !B$}
\DisplayProof
\hfill
\begin{tabular}{r}
\AxiomC{$!A \land !B$}
\RightLabel{\Elim{\land}}
\UnaryInfC{$!A$}
\DisplayProof
\\[3ex]
\AxiomC{$!A \land !B$}
\RightLabel{\Elim{\land}}
\UnaryInfC{$!B$}
\DisplayProof
\end{tabular}
\end{defish}

\subsection{Rules for $\lor$}

\begin{defish}
\begin{tabular}{r}
\AxiomC{$!A$}
\RightLabel{\Intro{\lor}}
\UnaryInfC{$!A \lor !B$}
\DisplayProof
\\[3ex]
\AxiomC{$!B$}
\RightLabel{\Intro{\lor}}
\UnaryInfC{$!A \lor !B$}
\DisplayProof
\end{tabular}
\hfill
\AxiomC{$!A \lor !B$}
\AxiomC{$\Discharge{!A}{n}$}
\DeduceC{$!C$}
\AxiomC{$\Discharge{!B}{n}$}
\DeduceC{$!C$}
\DischargeRule{\Elim{\lor}}{n}
\TrinaryInfC{$!C$}
\DisplayProof
\end{defish}

\subsection{Rules for $\lif$}

\begin{defish}
\AxiomC{$\Discharge{!A}{n}$}
\DeduceC{$!B$}
\DischargeRule{\Intro{\lif}}{n}
\UnaryInfC{$!A \lif !B$}
\DisplayProof
\hfill
\AxiomC{$!A \lif !B$}
\AxiomC{$!A$}
\RightLabel{\Elim{\lif}}
\BinaryInfC{$!B$}
\DisplayProof
\end{defish}

\subsection{Rules for $\lnot$}

\begin{defish}
\AxiomC{$\Discharge{!A}{n}$}
\noLine
\DeduceC{$\lfalse$}
\DischargeRule{\Intro{\lnot}}{n}
\UnaryInfC{$\lnot !A$}
\DisplayProof
\hfill
\AxiomC{$\lnot !A$}
\AxiomC{$!A$}
\RightLabel{\Elim{\lnot}}
\BinaryInfC{$\lfalse$}
\DisplayProof
\end{defish}

\subsection{Rules for $\lfalse$}

\begin{defish}
\AxiomC{$\lfalse$}
\RightLabel{\FalseInt}
\UnaryInfC{$!A$}
\DisplayProof
\hfill
\AxiomC{$\Discharge{\lnot !A}{n}$}
\DeduceC{$\lfalse$}
\DischargeRule{\FalseCl}{n}
\UnaryInfC{$!A$}
\DisplayProof
\end{defish}

Note that $\Intro{\lnot}$ and $\FalseCl$ are very similar: The
difference is that $\Intro{\lnot}$ derives a negated
!!{sentence}~$\lnot !A$ but $\FalseCl$ a positive !!{sentence}~$!A$.

Whenever a rule indicates that some assumption may be discharged, we
take this to be a permission, but not a requirement. E.g., in the
$\Intro{\lif}$ rule, we may discharge any number of assumptions of the
form~$!A$ in the !!{derivation} of the premise~$!B$, including zero.




  

\olfileid{fol}{ntd}{qrl}

\olsection{Quantifier Rules}

\subsection{Rules for $\lforall$}

\begin{defish}
\AxiomC{$!A(a)$}
\RightLabel{\Intro{\lforall}}
\UnaryInfC{$\lforall[x][\Atom{!A}{x}]$}
\DisplayProof
\hfill
\AxiomC{$\lforall[x][\Atom{!A}{x}]$}
\RightLabel{\Elim{\lforall}}
\UnaryInfC{$!A(t)$}
\DisplayProof
\end{defish}

In the rules for~$\lforall$, $t$ is a closed term (a term that does
not contain any variables), and $a$~is !!a{constant} which does not
occur in the conclusion~$\lforall[x][!A(x)]$, or in any assumption
which is !!{undischarged} in the !!{derivation} ending with the
premise~$!A(a)$. We call $a$ the \emph{eigenvariable} of the
\Intro{\lforall} inference.\footnote{We use the term ``eigenvariable''
even though $a$ in the above rule is a constant. This has historical
reasons.}

\subsection{Rules for $\lexists$}

\begin{defish}
\AxiomC{$\Atom{!A}{t}$}
\RightLabel{\Intro{\lexists}}
\UnaryInfC{$\lexists[x][\Atom{!A}{x}]$}
\DisplayProof
\hfill
\AxiomC{$\lexists[x][\Atom{!A}{x}]$}
\AxiomC{[$\Atom{!A}{a}$]$^n$}
\DeduceC{$!C$}
\DischargeRule{\Elim{\lexists}}{n}
\BinaryInfC{$!C$}
\DisplayProof
\end{defish}

Again, $t$ is a closed term, and $a$ is !!a{constant} which does not
occur in the premise $\lexists[x][!A(x)]$, in the conclusion~$!C$, or
any assumption which is !!{undischarged} in the !!{derivation}s ending
with the two premises (other than the assumptions $!A(a)$). We call
$a$ the \emph{eigenvariable} of the \Elim{\lexists} inference.

The condition that an eigenvariable neither occur in the premises nor
in any assumption that is !!{undischarged} in the !!{derivation}s
leading to the premises for the \Intro{\lforall} or \Elim{\lexists}
inference is called the \emph{eigenvariable condition}.

\begin{explain}
Recall the convention that when $!A$ is !!a{formula} with the
!!{variable}~$x$ free, we indicate this by writing~$!A(x)$. In the
same context, $!A(t)$ then is short for~$\Subst{!A}{t}{x}$. So we
could also write the $\Intro\lexists$ rule as:
\begin{prooftree}
\AxiomC{$\Subst{!A}{t}{x}$}
\RightLabel{\Intro{\lexists}}
\UnaryInfC{$\lexists[x][!A]$}
\end{prooftree}
Note that $t$ may already occur in~$!A$, e.g., $!A$~might
be~$\Atom{\Obj P}{t,x}$.  Thus, inferring $\lexists[x][\Atom{\Obj
P}{t,x}]$ from~$\Atom{\Obj P}{t,t}$ is a correct application
of~$\Intro\lexists$---you may ``replace'' one or more, and not
necessarily all, occurrences of~$t$ in the premise by the bound
!!{variable}~$x$. However, the eigenvariable conditions in
$\Intro\lforall$ and~$\Elim\lexists$ require that the !!{constant}~$a$
does not occur in~$!A$. So, you cannot correctly infer
$\lforall[x][\Atom{\Obj P}{a,x}]$ from $\Atom{\Obj P}{a,a}$
using~$\Intro\lforall$.
\end{explain}

\begin{explain}
In \Intro{\lexists} and \Elim{\lforall} there are no restrictions, and
the term~$t$ can be anything, so we do not have to worry about any
conditions. On the other hand, in the \Elim{\lexists} and
\Intro{\lforall} rules, the eigenvariable condition requires that the
!!{constant}~$a$ does not occur anywhere in the conclusion or in an
!!{undischarged} assumption. The condition is necessary to ensure that
the system is sound, i.e., only !!{derive}s !!{sentence}s from
!!{undischarged} assumptions from which they follow. Without this
condition, the following would be allowed:
\begin{prooftree}
  \AxiomC{$\lexists[x][!A(x)]$}
  \AxiomC{$\Discharge{!A(a)}{1}$}
  \RightLabel{*\Intro{\lforall}}
  \UnaryInfC{$\lforall[x][!A(x)]$}
  \RightLabel{\Elim{\lexists}}
  \BinaryInfC{$\lforall[x][!A(x)]$}
\end{prooftree}
However, $\lexists[x][!A(x)] \Entails/ \lforall[x][!A(x)]$.

As the elimination rules for quantifiers only allow substituting
closed terms for !!{variable}s, it follows that any !!{formula} that
can be derived from a set of !!{sentence}s is itself !!a{sentence}.
\end{explain}







\olfileid{fol}{ntd}{der}

\olsection{\usetoken{P}{derivation}}

\begin{explain}
We've said what an assumption is, and we've given the rules of
inference.  !!^{derivation}s in natural deduction are inductively
generated from these: each !!{derivation} either is an assumption
on its own, or consists of one, two, or three !!{derivation}s followed
by a correct inference.
\end{explain}

\begin{defn}[!!^{derivation}]
\Article{derivation} \emph{!!{derivation}} of !!a{sentence}~$!A$ from
assumptions~$\Gamma$ is a finite tree of !!{sentence}s satisfying the
following conditions:
\begin{enumerate}
\item The topmost !!{sentence}s of the tree are either in $\Gamma$ or
  are !!{discharged} by an inference in the tree.
\item The bottommost !!{sentence} of the tree is~$!A$.
\item Every !!{sentence} in the tree except the sentence~$!A$ at
  the bottom is a premise of a correct application of an inference
  rule whose conclusion stands directly below that !!{sentence} in the
  tree.
\end{enumerate}
We then say that $!A$ is the \emph{conclusion} of the !!{derivation}
and $\Gamma$ its !!{undischarged} assumptions.

If !!a{derivation} of $!A$ from~$\Gamma$ exists, we say that $!A$ is
\emph{!!{derivable}} from~$\Gamma$, or in symbols: $\Gamma \Proves
!A$. If there is !!a{derivation} of~$!A$ in which every assumption is
!!{discharged}, we write~$\Proves !A$.
\end{defn}

\begin{ex}
Every assumption on its own is !!a{derivation}. So, e.g., $!A$ by
itself is !!a{derivation}, and so is $!B$ by itself.  We can obtain a
new !!{derivation} from these by applying, say, the $\Intro{\land}$
rule,
\begin{prooftree}
\AxiomC{$!A$}
\AxiomC{$!B$}
\RightLabel{\Intro{\land}}
\BinaryInfC{$!A \land !B$}
\end{prooftree}
These rules are meant to be general: we can replace the $!A$ and~$!B$
in it with any !!{sentence}s, e.g., by $!C$ and~$!D$. Then the
conclusion would be $!C \land !D$, and so
\begin{prooftree}
\AxiomC{$!C$}
\AxiomC{$!D$}
\RightLabel{\Intro{\land}}
\BinaryInfC{$!C \land !D$}
\end{prooftree}
is a correct !!{derivation}. Of course, we can also switch the
assumptions, so that $!D$ plays the role of~$!A$ and $!C$ that
of~$!B$. Thus,
\begin{prooftree}
\AxiomC{$!D$}
\AxiomC{$!C$}
\RightLabel{\Intro{\land}}
\BinaryInfC{$!D \land !C$}
\end{prooftree}
is also a correct !!{derivation}.

We can now apply another rule, say, $\Intro{\lif}$, which allows us to
conclude a conditional and allows us to !!{discharge} any assumption
that is identical to the antecedent of that conditional. So both of
the following would be correct !!{derivation}s:
\begin{prooftree}
\AxiomC{$\Discharge{!C}{1}$}
\AxiomC{$!D$}
\RightLabel{\Intro{\land}}
\BinaryInfC{$!C \land !D$}
\DischargeRule{\Intro{\lif}}{1}
\UnaryInfC{$!C \lif (!C \land !D)$}
\DisplayProof\bottomAlignProof
\AxiomC{$!C$}
\AxiomC{$\Discharge{!D}{1}$}
\RightLabel{\Intro{\land}}
\BinaryInfC{$!C \land !D$}
\DischargeRule{\Intro{\lif}}{1}
\UnaryInfC{$!D \lif (!C \land !D)$}
\end{prooftree}
They show, respectively, that $!D \Proves !C \lif (!C \land !D)$ and
$!C \Proves !D \lif (!C \land !D)$.

Remember that discharging of assumptions is a permission, not a
requirement: we don't have to discharge the assumptions. In
particular, we can apply a rule even if the assumptions are not
present in the !!{derivation}. For instance, the following is legal, even
though there is no assumption~$!A$ to be !!{discharged}:
\begin{prooftree}
  \AxiomC{$!B$}
  \DischargeRule{\Intro{\lif}}{1}
  \UnaryInfC{$!A \lif !B$}
\end{prooftree}
\end{ex}





\olfileid{fol}{ntd}{pro}

\olsection{Examples of \usetoken{P}{derivation}}

\begin{ex}
Let's give !!a{derivation} of the !!{sentence} $(!A \land !B) \lif !A$.

We begin by writing the desired conclusion at the bottom of the 
!!{derivation}.
\begin{prooftree}
\AxiomC{}
\UnaryInfC{$(!A\land !B) \lif !A$}
\end{prooftree}

Next, we need to figure out what kind of inference could result in
!!a{sentence} of this form. The !!{main operator} of the
conclusion is $\lif$, so we'll try to arrive at the
conclusion using the \Intro{\lif} rule. It is best to write down
the assumptions involved and label the inference rules as you
progress, so it is easy to see whether all assumptions have been
!!{discharged} at the end of the proof.
\begin{prooftree}
\AxiomC{$\Discharge{!A \land !B}{1}$}
\DeduceC{$!A$}
\DischargeRule{\Intro{\lif}}{1} 
\UnaryInfC{$(!A\land !B) \lif !A$}
\end{prooftree}

We now need to fill in the steps from the assumption $!A \land !B$ to $!A$.
Since we only have one connective to deal with, $\land$, we must
use the $\land$ elim rule. This gives us the following proof:
\begin{prooftree}
\AxiomC{$\Discharge{!A \land !B}{1}$}
\RightLabel{\Elim{\land}}
\UnaryInfC{$!A$}
\DischargeRule{\Intro{\lif}}{1} 
\UnaryInfC{$(!A\land !B) \lif !A$}
\end{prooftree}
We now have a correct !!{derivation} of $(!A \land
!B) \lif !A$.
\end{ex}

\begin{ex}
Now let's give !!a{derivation} of $(\lnot !A \lor !B)
\lif (!A \lif !B)$.

We begin by writing the desired conclusion at the bottom of the 
!!{derivation}.
\begin{prooftree}
\AxiomC{}
\UnaryInfC{$(\lnot !A \lor !B) \lif (!A \lif !B)$}
\end{prooftree}
To find a logical rule that could give us this conclusion, we
look at the logical connectives in the conclusion: $\lnot$,
$\lor$, and $\lif$. We only care at the moment about the first
occurrence of $\lif$ because it is the !!{main operator} of the
!!{sentence} in the end-sequent, while $\lnot$, $\lor$ and the second
occurrence of $\lif$ are inside the scope of another connective, so we
will take care of those later. We therefore start with the
\Intro{\lif} rule.  A correct application must look like this:
\begin{prooftree}
\AxiomC{$\Discharge{\lnot !A \lor !B}{1}$}
\DeduceC{$!A \lif !B$}
\DischargeRule{\Intro{\lif}}{1}
\UnaryInfC{$(\lnot !A \lor !B) \lif (!A \lif !B)$}
\end{prooftree}
This leaves us with two possibilities to continue. Either we can
keep working from the bottom up and look for another application
of the \Intro{\lif} rule, or we can work from the top down and apply a
\Elim{\lor} rule. Let us apply the latter. We will use the assumption
$\lnot !A \lor !B$ as the leftmost premise of \Elim{\lor}.  For a valid
application of \Elim{\lor}, the other two premises must be identical
to the conclusion $!A \lif !B$, but each may be derived in turn from
another assumption, namely one of the two disjuncts of $\lnot !A \lor !B$.
So our !!{derivation} will look like this:
\begin{prooftree}
\AxiomC{$\Discharge{\lnot !A \lor !B}{1}$}
\AxiomC{$\Discharge{\lnot !A}{2}$}
\DeduceC{$!A \lif !B$}
\AxiomC{$\Discharge{!B}{2}$}
\DeduceC{$!A \lif !B$}
\DischargeRule{\Elim{\lor}}{2}
\TrinaryInfC{$!A \lif !B$}
\DischargeRule{\Intro{\lif}}{1} 
\UnaryInfC{$(\lnot !A \lor !B) \lif (!A \lif !B)$}
\end{prooftree}

In each of the two branches on the right, we want to !!{derive} $!A
\lif !B$, which is best done using \Intro{\lif}.
\begin{prooftree}
\AxiomC{$\Discharge{\lnot !A \lor !B}{1}$}
\AxiomC{$\Discharge{\lnot !A}{2}, \Discharge{!A}{3}$}
\DeduceC{$!B$}
\DischargeRule{\Intro{\lif}}{3}
\UnaryInfC{$!A \lif !B$}
\AxiomC{$\Discharge{!B}{2}, \Discharge{!A}{4}$}
\DeduceC{$!B$}
\DischargeRule{\Intro{\lif}}{4}
\UnaryInfC{$!A \lif !B$}
\DischargeRule{\Elim{\lor}}{2}
\TrinaryInfC{$!A \lif !B$}
\DischargeRule{\Intro{\lif}}{1} 
\UnaryInfC{$(\lnot !A \lor !B) \lif (!A \lif !B)$}
\end{prooftree}

For the two missing parts of the !!{derivation}, we need
!!{derivation}s of $!B$ from $\lnot !A$ and $!A$ in the middle, and
from $!A$ and $!B$ on the left.  Let's take the former first. $\lnot
!A$ and $!A$ are the two premises of \Elim{\lnot}:
\begin{prooftree}
\AxiomC{$\Discharge{\lnot !A}{2}$}
\AxiomC{$\Discharge{!A}{3}$}
\RightLabel{\Elim{\lnot}}
\BinaryInfC{$\lfalse$}
\DeduceC{$!B$}
\end{prooftree}
By using \FalseInt, we can obtain $!B$ as a conclusion and
complete the branch.
\begin{prooftree}
\AxiomC{$\Discharge{\lnot !A \lor !B}{1}$}
\AxiomC{$\Discharge{\lnot !A}{2}$}
\AxiomC{$\Discharge{!A}{3}$}
\RightLabel{\Intro{\lfalse}}
\BinaryInfC{$\lfalse$}
\RightLabel{\FalseInt}
\UnaryInfC{$!B$}
\DischargeRule{\Intro{\lif}}{3}
\UnaryInfC{$!A \lif !B$}
\AxiomC{$\Discharge{!B}{2}, \Discharge{!A}{4}$}
\DeduceC{$!B$}
\DischargeRule{\Intro{\lif}}{4}
\UnaryInfC{$!A \lif !B$}
\DischargeRule{\Elim{\lor}}{2}
\TrinaryInfC{$!A \lif !B$}
\DischargeRule{\Intro{\lif}}{1} 
\UnaryInfC{$(\lnot !A \lor !B) \lif (!A \lif !B)$}
\end{prooftree}

Let's now look at the rightmost branch.  Here it's important to
realize that the definition of !!{derivation} \emph{allows assumptions
  to be discharged} but \emph{does not require} them to be.  In other
words, if we can derive $!B$ from one of the assumptions $!A$ and $!B$
without using the other, that's ok.  And to !!{derive} $!B$ from~$!B$
is trivial: $!B$ by itself is such !!a{derivation}, and no inferences
are needed.  So we can simply delete the assumption~$!A$.
\begin{prooftree}
\AxiomC{$\Discharge{\lnot !A \lor !B}{1}$}
\AxiomC{$\Discharge{\lnot !A}{2}$}
\AxiomC{$\Discharge{!A}{3}$}
\RightLabel{\Elim{\lnot}}
\BinaryInfC{$\lfalse$}
\RightLabel{\FalseInt}
\UnaryInfC{$!B$}
\DischargeRule{\Intro{\lif}}{3}
\UnaryInfC{$!A \lif !B$}
\AxiomC{$\Discharge{!B}{2}$}
\RightLabel{\Intro{\lif}}
\UnaryInfC{$!A \lif !B$}
\DischargeRule{\Elim{\lor}}{2}
\TrinaryInfC{$!A \lif !B$}
\DischargeRule{\Intro{\lif}}{1} 
\UnaryInfC{$(\lnot !A \lor !B) \lif (!A \lif !B)$}
\end{prooftree}
Note that in the finished !!{derivation}, the rightmost \Intro{\lif}
inference does not actually discharge any assumptions.
\end{ex}

\begin{ex}
So far we have not needed the \FalseCl{} rule. It is special in that
it allows us to discharge an assumption that isn't a sub-!!{formula} of
the conclusion of the rule.  It is closely related to the \FalseInt{}
rule. In fact, the \FalseInt{} rule is a special case of the
\FalseCl{} rule---there is a logic called ``intuitionistic logic'' in
which only \FalseInt{} is allowed.  The \FalseCl{} rule is a last
resort when nothing else works.  For instance, suppose we want to
!!{derive} $!A \lor \lnot !A$. Our usual strategy would be to attempt
to !!{derive} $!A \lor \lnot !A$ using $\Intro{\lor}$. But this would
require us to !!{derive} either $!A$ or $\lnot !A$ from no
assumptions, and this can't be done. \FalseCl{} to the rescue!
\begin{prooftree}
  \AxiomC{$\Discharge{\lnot(!A \lor \lnot !A)}{1}$}
  \DeduceC{$\lfalse$}
  \DischargeRule{\FalseCl}{1}
  \UnaryInfC{$!A \lor \lnot !A$}
\end{prooftree}
Now we're looking for !!a{derivation} of $\lfalse$ from $\lnot(!A \lor
\lnot !A)$. Since $\lfalse$ is the conclusion of $\Elim{\lnot}$ we
might try that:
\begin{prooftree}
  \AxiomC{$\Discharge{\lnot(!A \lor \lnot !A)}{1}$}
  \DeduceC{$\lnot !A$}
  \AxiomC{$\Discharge{\lnot(!A \lor \lnot !A)}{1}$}
  \DeduceC{$!A$}
  \RightLabel{\Elim{\lnot}}
  \BinaryInfC{$\lfalse$}
  \DischargeRule{\FalseCl}{1}
  \UnaryInfC{$!A \lor \lnot !A$}
\end{prooftree}
Our strategy for finding !!a{derivation} of~$\lnot !A$ calls for an
application of~$\Intro{\lnot}$:
\begin{prooftree}
  \AxiomC{$\Discharge{\lnot(!A \lor \lnot !A)}{1}, \Discharge{!A}{2}$}
  \DeduceC{$\lfalse$}
  \DischargeRule{\Intro{\lnot}}{2}
  \UnaryInfC{$\lnot !A$}
  \AxiomC{$\Discharge{\lnot(!A \lor \lnot !A)}{1}$}
  \DeduceC{$!A$}
  \RightLabel{\Elim{\lnot}}
  \BinaryInfC{$\lfalse$}
  \DischargeRule{\FalseCl}{1}
  \UnaryInfC{$!A \lor \lnot !A$}
\end{prooftree}
Here, we can get $\lfalse$ easily by applying $\Elim{\lnot}$ to the
assumption $\lnot(!A \lor \lnot !A)$ and $!A \lor \lnot !A$ which
follows from our new assumption $!A$ by~$\Intro{\lor}$:
\begin{prooftree}
  \AxiomC{$\Discharge{\lnot(!A \lor \lnot !A)}{1}$}
  \AxiomC{$\Discharge{!A}{2}$}
  \RightLabel{\Intro{\lor}}
  \UnaryInfC{$!A \lor \lnot !A$}
  \RightLabel{\Elim{\lnot}}
  \BinaryInfC{$\lfalse$}
  \DischargeRule{\Intro{\lnot}}{2}
  \UnaryInfC{$\lnot !A$}
  \AxiomC{$\Discharge{\lnot(!A \lor \lnot !A)}{1}$}
  \DeduceC{$!A$}
  \RightLabel{\Elim{\lnot}}
  \BinaryInfC{$\lfalse$}
  \DischargeRule{\FalseCl}{1}
  \UnaryInfC{$!A \lor \lnot !A$}
\end{prooftree}
On the right side we use the same strategy, except we get $!A$ by~\FalseCl:
\begin{prooftree}
  \AxiomC{$\Discharge{\lnot(!A \lor \lnot !A)}{1}$}
  \AxiomC{$\Discharge{!A}{2}$}
  \RightLabel{\Intro{\lor}}
  \UnaryInfC{$!A \lor \lnot !A$}
  \RightLabel{\Elim{\lnot}}
  \BinaryInfC{$\lfalse$}
  \DischargeRule{\Intro{\lnot}}{2}
  \UnaryInfC{$\lnot !A$}
  \AxiomC{$\Discharge{\lnot(!A \lor \lnot !A)}{1}$}
  \AxiomC{$\Discharge{\lnot !A}{3}$}
  \RightLabel{\Intro{\lor}}
  \UnaryInfC{$!A \lor \lnot !A$}
  \RightLabel{\Elim{\lnot}}
  \BinaryInfC{$\lfalse$}
  \DischargeRule{\FalseCl}{3}
  \UnaryInfC{$!A$}
  \RightLabel{\Elim{\lnot}}
  \BinaryInfC{$\lfalse$}
  \DischargeRule{\FalseCl}{1}
  \UnaryInfC{$!A \lor \lnot !A$}
\end{prooftree}
\end{ex}

\begin{prob}
Give !!{derivation}s that show the following:
\begin{enumerate}
\item $!A \land (!B \land !C) \Proves (!A \land !B) \land !C$.
\item $!A \lor (!B \lor !C) \Proves (!A \lor !B) \lor !C$.
\item $!A \lif (!B \lif !C) \Proves !B \lif (!A \lif !C)$.
\item $!A \Proves \lnot\lnot !A$.
\end{enumerate}
\end{prob}

\begin{prob}
Give !!{derivation}s that show the following:
\begin{enumerate}
\item $(!A \lor !B) \lif !C \Proves !A \lif !C$.
\item $(!A \lif !C) \land (!B \lif !C) \Proves (!A \lor !B) \lif !C$.
\item $\Proves \lnot(!A \land \lnot !A)$.
\item $!B \lif !A \Proves \lnot !A \lif \lnot !B$.
\item $\Proves (!A \lif \lnot !A) \lif \lnot !A$.
\item $\Proves \lnot(!A \lif !B) \lif \lnot !B$.
\item $!A \lif !C \Proves \lnot (!A \land \lnot !C)$.
\item $!A \land \lnot !C \Proves \lnot (!A \lif !C)$.
\item $!A \lor !B, \lnot !B \Proves !A$.
\item $\lnot !A \lor \lnot !B \Proves \lnot(!A \land !B)$.
\item $\Proves (\lnot !A \land \lnot !B) \lif\lnot(!A \lor !B)$.
\item $\Proves \lnot(!A \lor !B) \lif (\lnot !A \land \lnot !B)$.
\end{enumerate}
\end{prob}

\begin{prob}
Give !!{derivation}s that show the following:
\begin{enumerate}
\item $\lnot(!A \lif !B) \Proves !A$.
\item $\lnot(!A \land !B) \Proves \lnot !A \lor \lnot !B$.
\item $!A \lif !B \Proves \lnot !A \lor !B$.
\item $\Proves \lnot \lnot !A \lif !A$.
\item $!A \lif !B, \lnot !A \lif !B \Proves !B$.
\item $(!A \land !B) \lif !C \Proves (!A \lif !C) \lor (!B \lif !C)$.
\item $(!A \lif !B) \lif !A \Proves !A$.
\item $\Proves (!A \lif !B) \lor (!B \lif !C)$.
\end{enumerate}
(These all require the $\FalseCl$~rule.)
\end{prob}




  

\olfileid{fol}{ntd}{prq}

\olsection{\usetoken{P}{derivation} with Quantifiers}

\begin{ex}
When dealing with quantifiers, we have to make sure not to violate the
eigenvariable condition, and sometimes this requires us to play around
with the order of carrying out certain inferences. In general, it
helps to try and take care of rules subject to the eigenvariable
condition first (they will be lower down in the finished proof).

Let's see how we'd give !!a{derivation} of the !!{formula}
$\lexists[x][\lnot !A(x)] \lif \lnot \lforall[x][!A(x)]$.
Starting as usual, we write
\begin{prooftree}
\AxiomC{}
\UnaryInfC{$\lexists[x][\lnot !A(x)]\lif \lnot \lforall[x][!A(x)]$}
\end{prooftree}
We start by writing down what it would take to justify that last step
using the \Intro{\lif} rule.
\begin{prooftree}
\AxiomC{$\Discharge{\lexists[x][\lnot !A(x)]}{1}$}
\DeduceC{$\lnot \lforall[x][!A(x)]$}
\DischargeRule{\Intro{\lif}}{1}
\UnaryInfC{$\lexists[x][\lnot !A(x)]\lif \lnot \lforall[x][!A(x)]$}
\end{prooftree}
Since there is no obvious rule to apply to $\lnot \lforall[x][!A(x)]$,
we will proceed by setting up the !!{derivation} so we can use the
\Elim{\lexists} rule. Here we must pay attention to the eigenvariable
condition, and choose a constant that does not appear in
$\lexists[x][!A(x)]$ or any assumptions that it depends on.
(Since no !!{constant}s appear, however, any choice will do fine.)
\begin{prooftree}
\AxiomC{$\Discharge{\lexists[x][\lnot !A(x)]}{1}$}
\AxiomC{$\Discharge{\lnot !A(a)}{2}$}
\DeduceC{$\lnot \lforall[x][!A(x)]$}
\DischargeRule{\Elim{\lexists}}{2}
\BinaryInfC{$\lnot \lforall[x][!A(x)]$}
\DischargeRule{\Intro{\lif}}{1}
\UnaryInfC{$\lexists[x][\lnot !A(x)] \lif \lnot \lforall[x][!A(x)]$}
\end{prooftree}
In order to derive $\lnot \lforall[x][!A(x)]$, we will attempt to use
the \Intro{\lnot} rule: this requires that we derive a contradiction,
possibly using $\lforall[x][!A(x)]$ as an additional assumption. Of
course, this contradiction may involve the assumption $\lnot !A(a)$
which will be discharged by the \Elim{\lexists} inference. We can set it
up as follows:
\begin{prooftree}
\AxiomC{$\Discharge{\lexists[x][\lnot !A(x)]}{1}$}
\AxiomC{$\Discharge{\lnot !A(a)}{2}, \Discharge{\lforall[x][!A(x)]}{3}$}
\DeduceC{$\lfalse$}
\DischargeRule{\Intro{\lnot}}{3}
\UnaryInfC{$\lnot \lforall[x][!A(x)]$}
\DischargeRule{\Elim{\lexists}}{2}
\BinaryInfC{$\lnot \lforall[x][!A(x)]$}
\DischargeRule{\Intro{\lif}}{1}
\UnaryInfC{$\lexists[x][\lnot !A(x)]\lif \lnot \lforall[x][!A(x)]$}
\end{prooftree}
It looks like we are close to getting a contradiction. The easiest
rule to apply is the \Elim{\lforall}, which has no eigenvariable
conditions. Since we can use any term we want to replace the
universally quantified~$x$, it makes the most sense to continue
using~$a$ so we can reach a contradiction.
\begin{prooftree}
\AxiomC{$\Discharge{\lexists[x][\lnot !A(x)]}{1}$}
\AxiomC{$\Discharge{\lnot !A(a)}{2}$}
\AxiomC{$\Discharge{\lforall[x][!A(x)]}{3}$}
\RightLabel{\Elim{\lforall}}
\UnaryInfC{$!A(a)$}
\RightLabel{\Elim{\lnot}}
\BinaryInfC{$\lfalse$}
\DischargeRule{\Intro{\lnot}}{3}
\UnaryInfC{$\lnot \lforall[x][!A(x)]$}
\DischargeRule{\Elim{\lexists}}{2}
\BinaryInfC{$\lnot \lforall[x][!A(x)]$}
\DischargeRule{\Intro{\lif}}{1}
\UnaryInfC{$\lexists[x][\lnot !A(x)]\lif \lnot \lforall[x][!A(x)]$}
\end{prooftree}

It is important, especially when dealing with quantifiers, to double
check at this point that the eigenvariable condition has not been
violated. Since the only rule we applied that is subject to the
eigenvariable condition was \Elim{\exists}, and the eigenvariable~$a$
does not occur in any assumptions it depends on, this is a
correct !!{derivation}.
\end{ex}

\begin{ex}
Sometimes we may derive !!a{formula} from other !!{formula}s.
In these cases, we may have undischarged assumptions. It is 
important to keep track of our assumptions as well
as the end goal.

Let's see how we'd give !!a{derivation} of the !!{formula}
$\lexists[x][!C(x,b)]$ from the assumptions $\lexists[x][(!A(x) 
\land !B(x))]$ and $\lforall[x][(!B(x) \lif !C(x,b))]$.
Starting as usual, we write the conclusion at the
bottom.
\begin{prooftree}
\AxiomC{}
\UnaryInfC{$\lexists[x][!C(x,b)]$}
\end{prooftree}

We have two premises to work with. To use the first, i.e., try to find
!!a{derivation} of $\lexists[x][!C(x, b)]$ from $\lexists[x][(!A(x)
  \land !B(x))]$ we would use the \Elim{\lexists} rule. Since it has
an eigenvariable condition, we will apply that rule first. We get the
following:
\begin{prooftree}
\AxiomC{$\lexists[x][(!A(x) \land !B(x))]$}
\AxiomC{$\Discharge{!A(a) \land !B(a)}{1}$}
\DeduceC{$\lexists[x][!C(x,b)]$}
\DischargeRule{\Elim{\lexists}}{1}
\BinaryInfC{$\lexists[x][!C(x,b)]$}
\end{prooftree}
The two assumptions we are working with share~$!B$.  It may be useful
at this point to apply \Elim{\land} to separate out~$!B(a)$.
\begin{prooftree}
\AxiomC{$\lexists[x][(!A(x) \land !B(x)])$}
\AxiomC{$\Discharge{!A(a) 
\land !B(a)}{1}$}
\RightLabel{\Elim{\land}}
\UnaryInfC{$!B(a)$}
\DeduceC{$\lexists[x][!C(x,b)]$}
\DischargeRule{\Elim{\lexists}}{1}
\BinaryInfC{$\lexists[x][!C(x,b)]$}
\end{prooftree}

The second assumption we have to work with is~$\lforall[x][(!B(x) \lif
  !C(x,b))]$. Since there is no eigenvariable condition we can
instantiate $x$ with the !!{constant}~$a$ using \Elim{\lforall} to get
$!B(a) \lif !C(a, b)$.  We now have both $!B(a) \lif !C(a,b)$ and
$!B(a)$. Our next move should be a straightforward application of the
\Elim{\lif} rule.
\begin{prooftree}
\AxiomC{$\lexists[x][(!A(x) \land !B(x))]$}
\AxiomC{$\lforall[x][(!B(x) \lif !C(x,b))]$}
\RightLabel{\Elim{\lforall}}
\UnaryInfC{$!B(a) \lif !C(a,b)$}
\AxiomC{$\Discharge{!A(a) 
\land !B(a)}{1}$}
\RightLabel{\Elim{\land}}
\UnaryInfC{$!B(a)$}
\RightLabel{\Elim{\lif}}
\BinaryInfC{$!C(a,b)$}
\DeduceC{$\lexists[x][!C(x,b)]$}
\DischargeRule{\Elim{\lexists}}{1}
\insertBetweenHyps{\hspace{-5em}}
\BinaryInfC{$\lexists[x][!C(x,b)]$}
\end{prooftree}
We are so close!{} One application of \Intro{\lexists} and we
have reached our goal.
\begin{prooftree}
\AxiomC{$\lexists[x][(!A(x) \land !B(x))]$}
\AxiomC{$\lforall[x][(!B(x) \lif !C(x,b))]$}
\RightLabel{\Elim{\lforall}}
\UnaryInfC{$!B(a) \lif !C(a,b)$}
\AxiomC{$\Discharge{!A(a) 
\land !B(a)}{1}$}
\RightLabel{\Elim{\land}}
\UnaryInfC{$!B(a)$}
\RightLabel{\Elim{\lif}}
\BinaryInfC{$!C(a,b)$}
\RightLabel{\Intro{\lexists}}
\UnaryInfC{$\lexists[x][!C(x,b)]$}
\DischargeRule{\Elim{\lexists}}{1}
\insertBetweenHyps{\hspace{-5em}}
\BinaryInfC{$\lexists[x][!C(x,b)]$}
\end{prooftree}
Since we ensured at each step that the eigenvariable
conditions were not violated, we can be confident that this
is a correct !!{derivation}.
\end{ex}

\begin{ex}
Give !!a{derivation} of the !!{formula}
$\lnot\lforall[x][!A(x)]$ from the assumptions $\lforall[x][!A(x)] 
\lif \lexists[y][!B(y)]$ and $\lnot\lexists[y][!B(y)]$.
Starting as usual, we write the target !!{formula} at the bottom.
\begin{prooftree}
\AxiomC{}
\UnaryInfC{$\lnot\lforall[x][!A(x)]$}
\end{prooftree}
The last line of the !!{derivation} is a negation, so let's try using
\Intro{\lnot}. This will require that we figure out how to !!{derive}
a contradiction.
\begin{prooftree}
\AxiomC{$\Discharge{\lforall[x][!A(x)]}{1}$}
\DeduceC{$\lfalse$}
\DischargeRule{\Intro{\lnot}}{1}
\UnaryInfC{$\lnot\lforall[x][!A(x)]$}
\end{prooftree}
So far so good. We can use \Elim{\lforall} but it's not obvious
if that will help us get to our goal. Instead, let's use one of our 
assumptions. $\lforall[x][!A(x)] \lif \lexists[y][!B(y)]$ together
with $\lforall[x][!A(x)]$ will allow us to use the \Elim{\lif} rule.
\begin{prooftree}
\AxiomC{$\lforall[x][!A(x)] \lif \lexists[y][!B(y)]$}
\AxiomC{$\Discharge{\lforall[x][!A(x)]}{1}$}
\RightLabel{\Elim{\lif}}
\BinaryInfC{$\lexists[y][!B(y)]$}
\DeduceC{$\lfalse$}
\DischargeRule{\Intro{\lnot}}{1}
\UnaryInfC{$\lnot\lforall[x][!A(x)]$}
\end{prooftree}
We now have one final assumption to work with,
and it looks like this will help us reach a contradiction
by using \Elim{\lnot}.
\begin{prooftree}
\AxiomC{$\lnot\lexists[y][!B(y)]$}
\AxiomC{$\lforall[x][!A(x)] \lif \lexists[y][!B(y)]$}
\AxiomC{$\Discharge{\lforall[x][!A(x)]}{1}$}
\RightLabel{\Elim{\lif}}
\BinaryInfC{$\lexists[y][!B(y)]$}
\RightLabel{\Elim{\lnot}}
\BinaryInfC{$\lfalse$}
\DischargeRule{\Intro{\lnot}}{1}
\UnaryInfC{$\lnot\lforall[x][!A(x)]$}
\end{prooftree}
\end{ex}

\begin{prob}
Give !!{derivation}s that show the following:
\begin{enumerate}
\item $\Proves (\lforall[x][!A(x)] \land \lforall[y][!B(y)]) \lif
\lforall[z][(!A(z) \land !B(z))]$.
\item $\Proves (\lexists[x][!A(x)] \lor \lexists[y][!B(y)]) \lif
\lexists[z][(!A(z) \lor !B(z))]$.
\item $\lforall[x][(!A(x) \lif !B)] \Proves \lexists[y][!A(y)] \lif !B$.
\item $\lforall[x][\lnot !A(x)] \Proves \lnot\lexists[x][!A(x)]$.
\item $\Proves \lnot\lexists[x][!A(x)] \lif \lforall[x][\lnot !A(x)]$.
\item $\Proves \lnot\lexists[x][\lforall[y][((!A(x,y) \lif \lnot
!A(y,y)) \land (\lnot !A(y,y) \lif !A(x,y)))]]$.
\end{enumerate}
\end{prob}

\begin{prob}
Give !!{derivation}s that show the following:
\begin{enumerate}
\item $\Proves \lnot\lforall[x][!A(x)] \lif \lexists[x][\lnot!A(x)]$.
\item $(\lforall[x][!A(x)] \lif !B) \Proves \lexists[y][(!A(y) \lif !B)]$.
\item $\Proves \lexists[x][(!A(x) \lif \lforall[y][!A(y)])]$.
\end{enumerate}
(These all require the $\FalseCl$~rule.)
\end{prob}






\olfileid{fol}{ntd}{ptn}

\olsection{Proof-Theoretic Notions}

\begin{editorial}
This section collects the definitions the provability relation
and consistency for natural deduction.
\end{editorial}

\begin{explain}
Just as we've defined a number of important semantic notions
(validity, entailment, satisfiability), we now define corresponding
\emph{proof-theoretic notions}.  These are not defined by appeal to
satisfaction of !!{sentence}s in !!{structure}s, but by appeal to the
!!{derivability} or !!{nonderivability} of certain !!{sentence}s from
others.  It was an important discovery that these notions coincide.
That they do is the content of the \emph{soundness} and
\emph{completeness theorems}.
\end{explain}


\begin{defn}[Theorems]
A !!{sentence}~$!A$ is a \emph{theorem} if there is !!a{derivation}
of~$!A$ in natural deduction in which all assumptions are
!!{discharged}.  We write $\Proves !A$ if $!A$ is a theorem and
$\Proves/ !A$ if it is not.
\end{defn}

\begin{defn}[!!^{derivability}]
!!^a{sentence} $!A$ is \emph{!!{derivable} from} a set of
!!{sentence}s~$\Gamma$, $\Gamma \Proves !A$, if there is a
!!{derivation} with conclusion~$!A$ and in which every assumption
is either !!{discharged} or is in~$\Gamma$. If $!A$ is not
!!{derivable} from $\Gamma$ we write $\Gamma \Proves/ !A$.
\end{defn}

\begin{defn}[Consistency]
A set of !!{sentence}s~$\Gamma$ is \emph{inconsistent} iff $\Gamma
\Proves \lfalse$.  If $\Gamma$ is not inconsistent, i.e., if
$\Gamma \Proves/ \lfalse$, we say it is \emph{consistent}.
\end{defn}

\begin{prop}[Reflexivity]
\ollabel{prop:reflexivity}
If $!A \in \Gamma$, then $\Gamma \Proves !A$.
\end{prop}

\begin{proof}
The assumption $!A$ by itself is !!a{derivation} of~$!A$ where every
!!{undischarged} assumption (i.e., $!A$) is in~$\Gamma$.
\end{proof}
  

\begin{prop}[Monotonicity]
\ollabel{prop:monotonicity}
If $\Gamma \subseteq \Delta$ and $\Gamma \Proves !A$, then $\Delta
\Proves !A$.
\end{prop}

\begin{proof}
Any !!{derivation} of $!A$ from $\Gamma$ is also !!a{derivation} of
$!A$ from~$\Delta$.
\end{proof}

\begin{prop}[Transitivity]
\ollabel{prop:transitivity}
If $\Gamma \Proves !A$ and $\{!A\} \cup \Delta \Proves
!B$, then $\Gamma \cup \Delta \Proves !B$.
\end{prop}

\begin{proof}
If $\Gamma \Proves !A$, there is !!a{derivation}~$\delta_0$ of~$!A$
with all !!{undischarged} assumptions in~$\Gamma$.  If $\{!A\} \cup
\Delta \Proves !B$, then there is !!a{derivation}~$\delta_1$ of~$!B$
with all !!{undischarged} assumptions in~$\{!A\} \cup \Delta$.
Now consider:
\begin{prooftree}
  \AxiomC{$\Delta, \Discharge{!A}{1}$}
  \RightLabel{$\delta_1$}
  \DeduceC{$!B$}
  \DischargeRule{\Intro{\lif}}{1}
  \UnaryInfC{$!A \lif !B$}
  \AxiomC{$\Gamma$}
  \RightLabel{$\delta_0$}
  \DeduceC{$!A$}
  \RightLabel{\Elim{\lif}}
  \BinaryInfC{$!B$}
\end{prooftree}
The !!{undischarged} assumptions are now all among $\Gamma \cup
\Delta$, so this shows $\Gamma \cup \Delta \Proves !B$.
\end{proof}

When $\Gamma = \{!A_1, !A_2, \ldots, !A_k\}$ is a finite set we may use the simplified notation $!A_1,!A_2,\ldots,!A_k \Proves !B$ for $\Gamma \Proves !B$, in particular $!A \Proves !B$ means that $\{!A\} \Proves !B$.

Note that if $\Gamma \Proves !A$ and $!A
\Proves !B$, then $\Gamma \Proves !B$. It follows also that if $!A_1,
\dots, !A_n \Proves !B$ and $\Gamma \Proves !A_i$ for each~$i$, then
$\Gamma \Proves !B$.

\begin{prop}
\ollabel{prop:incons}
The following are equivalent. 
    \begin{enumerate}
      \item \( \Gamma \) is inconsistent.
      \item \( \Gamma \Proves {!A} \) for every !!{sentence}~\( {!A} \).
      \item \( \Gamma \Proves {!A} \) and \( \Gamma \Proves \lnot {!A} \) for some !!{sentence}~\( {!A} \).
    \end{enumerate}
\end{prop}

\begin{proof}
Exercise.
\end{proof}


\begin{prob}
Prove \olref[fol][ntd][ptn]{prop:incons}
\end{prob}




\begin{prop}[Compactness]
  \ollabel{prop:proves-compact}
  \begin{enumerate}
  \item If $\Gamma \Proves !A$ then there is a finite subset $\Gamma_0
    \subseteq \Gamma$ such that $\Gamma_0 \Proves !A$.
  \item If every finite subset of~$\Gamma$ is
    consistent, then $\Gamma$ is consistent.
  \end{enumerate}
\end{prop}

\begin{proof}
  \begin{enumerate}
    \item If $\Gamma \Proves !A$, then there is
      !!a{derivation}~$\delta$ of~$!A$ from~$\Gamma$. Let $\Gamma_0$
      be the set of !!{undischarged} assumptions of~$\delta$.  Since
      any !!{derivation} is finite, $\Gamma_0$ can only contain
      finitely many !!{sentence}s.  So, $\delta$ is !!a{derivation}
      of~$!A$ from a finite~$\Gamma_0 \subseteq \Gamma$.
    \item This is the contrapositive of (1) for the special case $!A
      \ident \lfalse$.
  \end{enumerate}
\end{proof}





\olfileid{fol}{ntd}{prv}
      
\olsection{\usetoken{S}{derivability} and Consistency}

We will now establish a number of properties of the !!{derivability}
relation.  They are independently interesting, but each will play a
role in the proof of the completeness theorem.

\begin{prop}\ollabel{prop:provability-contr}
  If $\Gamma \Proves !A$ and $\Gamma \cup \{!A\}$ is inconsistent,
  then $\Gamma$ is inconsistent.
\end{prop}

\begin{proof}
Let the !!{derivation} of~$!A$ from~$\Gamma$ be~$\delta_1$ and the
!!{derivation} of~$\lfalse$ from $\Gamma \cup \{!A\}$
be~$\delta_2$. We can then !!{derive}:
\begin{prooftree}
\AxiomC{$\Gamma, \Discharge{!A}{1}$}
\RightLabel{$\delta_2$}
\DeduceC{$\lfalse$}
\DischargeRule{\Intro{\lnot}}{1}
\UnaryInfC{$\lnot !A$}
\AxiomC{$\Gamma$}
\RightLabel{$\delta_1$}
\DeduceC{$!A$}
\RightLabel{\Elim{\lnot}}
\BinaryInfC{$\lfalse$}
\end{prooftree}
In the new !!{derivation}, the assumption~$!A$ is !!{discharged}, so it is
!!a{derivation} from~$\Gamma$.
\end{proof}

\begin{prop}
\ollabel{prop:prov-incons}
$\Gamma \Proves !A$ iff $\Gamma \cup \{\lnot !A\}$ is inconsistent.
\end{prop}

\begin{proof}
First suppose $\Gamma \Proves !A$, i.e., there is
!!a{derivation}~$\delta_0$ of~$!A$ from !!{undischarged}
assumptions~$\Gamma$. We obtain !!a{derivation} of $\lfalse$ from
$\Gamma \cup \{\lnot !A\}$ as follows:
\begin{prooftree}
  \AxiomC{$\lnot !A$}
  \AxiomC{$\Gamma$}
  \RightLabel{$\delta_0$}
  \DeduceC{$!A$}
  \RightLabel{\Elim{\lnot}}
  \BinaryInfC{$\lfalse$}
\end{prooftree}

Now assume $\Gamma \cup \{\lnot !A\}$ is inconsistent, and let
$\delta_1$ be the corresponding !!{derivation} of~$\lfalse$ from
!!{undischarged} assumptions in~$\Gamma \cup \{\lnot !A\}$. We obtain
!!a{derivation} of~$!A$ from~$\Gamma$ alone by using~$\FalseCl$:
\begin{prooftree}
  \AxiomC{$\Gamma, \Discharge{\lnot !A}{1}$}
  \RightLabel{$\delta_1$}
  \DeduceC{$\lfalse$}
  \RightLabel{\FalseCl}
  \DischargeRule{\FalseCl}{1}
  \UnaryInfC{$!A$}
\end{prooftree}
\end{proof}

\begin{prob}
Prove that $\Gamma \Proves \lnot !A$ iff $\Gamma \cup \{!A\}$ is
inconsistent.
\end{prob}

\begin{prop}\ollabel{prop:explicit-inc}
  If $\Gamma \Proves !A$ and $\lnot !A \in \Gamma$, then $\Gamma$ is
  inconsistent.
\end{prop}

\begin{proof}
  Suppose $\Gamma \Proves !A$ and $\lnot !A \in \Gamma$.  Then there
  is !!a{derivation}~$\delta$ of~$!A$ from~$\Gamma$. Consider this
    simple application of the $\Elim{\lnot}$ rule:
  \begin{prooftree}
    \AxiomC{$\lnot !A$}
    \AxiomC{$\Gamma$}
    \RightLabel{$\delta$}
    \DeduceC{$!A$}
    \RightLabel{\Elim{\lnot}}
    \BinaryInfC{$\lfalse$}
  \end{prooftree}
  Since $\lnot !A \in \Gamma$, all !!{undischarged} assumptions are
  in~$\Gamma$, this shows that $\Gamma \Proves \lfalse$.
\end{proof}

\begin{prop}\ollabel{prop:provability-exhaustive}
  If $\Gamma \cup \{!A\}$ and $\Gamma \cup \{\lnot !A\}$ are both
  inconsistent, then $\Gamma$ is inconsistent.
\end{prop}

\begin{proof}
There are !!{derivation}s $\delta_1$ and $\delta_2$ of~$\lfalse$ from
  $\Gamma \cup \{ !A \}$ and $\lfalse$ from $\Gamma \cup \{ \lnot !A
  \}$, respectively. We can then !!{derive}
\begin{prooftree}
\AxiomC{$\Gamma, \Discharge{\lnot !A}{2}$}
\RightLabel{$\delta_2$}
\DeduceC{$\lfalse$}
\DischargeRule{\Intro{\lnot}}{2}
\UnaryInfC{$\lnot \lnot !A$}
\AxiomC{$\Gamma, \Discharge{!A}{1}$}
\RightLabel{$\delta_1$}
\DeduceC{$\lfalse$}
\DischargeRule{\Intro{\lnot}}{1}
\UnaryInfC{$\lnot !A$}
\RightLabel{\Elim{\lnot}}
\BinaryInfC{$\lfalse$}
\end{prooftree}
Since the assumptions $!A$ and $\lnot !A$ are !!{discharged}, this is
!!a{derivation} of~$\lfalse$ from~$\Gamma$ alone. Hence $\Gamma$ is
inconsistent.
\end{proof}





\olfileid{fol}{ntd}{ppr}

\olsection{\usetoken{S}{derivability} and the Propositional Connectives}

\begin{explain}
  We establish that the !!{derivability} relation~$\Proves$ of natural
 deduction is strong enough to establish some basic facts
  involving the propositional connectives, such as that $!A \land !B
  \Proves !A$ and $!A, !A \lif !B \Proves !B$ (modus ponens). These
  facts are needed for the proof of the completeness theorem.
\end{explain}

\begin{prop}\ollabel{prop:provability-land}
  \begin{enumerate}
  \item \ollabel{prop:provability-land-left} Both $!A \land !B \Proves
    !A$ and $!A \land !B \Proves !B$
  \item \ollabel{prop:provability-land-right} $!A, !B \Proves !A \land !B$.
  \end{enumerate}
\end{prop}

\begin{proof}
  \begin{enumerate}
  \item We can !!{derive} both
    \begin{prooftree}
      \AxiomC{$!A \land !B$}
      \RightLabel{\Elim{\land}}
      \UnaryInfC{$!A$}
      \DisplayProof\qquad\bottomAlignProof
      \AxiomC{$!A \land !B$}
      \RightLabel{\Elim{\land}}
      \UnaryInfC{$!B$}
    \end{prooftree}
  \item We can !!{derive}:
    \begin{prooftree}
      \AxiomC{$!A$}
      \AxiomC{$!B$}
      \RightLabel{\Intro{\land}}
      \BinaryInfC{$!A \land !B$}
    \end{prooftree}
  \end{enumerate}
\end{proof}


\begin{prop}\ollabel{prop:provability-lor}
  \begin{enumerate}
  \item $!A \lor !B, \lnot !A, \lnot !B$ is inconsistent.
  \item Both $!A \Proves !A \lor !B$ and $!B \Proves !A \lor !B$.
  \end{enumerate}
\end{prop}

\begin{proof}
  \begin{enumerate}
  \item Consider the following !!{derivation}:
    \begin{prooftree}
      \AxiomC{$!A \lor !B$}
      \AxiomC{$\lnot !A$}
      \AxiomC{$\Discharge{!A}{1}$}
      \RightLabel{\Elim{\lnot}}
      \BinaryInfC{$\lfalse$}
      \AxiomC{$\lnot !B$}
      \AxiomC{$\Discharge{!B}{1}$}
      \RightLabel{\Elim{\lnot}}
      \BinaryInfC{$\lfalse$}
      \DischargeRule{\Elim{\lor}}{1}
      \TrinaryInfC{$\lfalse$}
    \end{prooftree}
    This is !!a{derivation} of~$\lfalse$ from !!{undischarged}
    assumptions $!A \lor !B$, $\lnot !A$, and $\lnot !B$.
  \item We can !!{derive} both
    \begin{prooftree}
      \AxiomC{$!A$}
      \RightLabel{\Intro{\lor}}
      \UnaryInfC{$!A \lor !B$}
      \DisplayProof\qquad\bottomAlignProof
      \AxiomC{$!B$}
      \RightLabel{\Intro{\lor}}
      \UnaryInfC{$!A \lor !B$}
    \end{prooftree}
  \end{enumerate}
\end{proof}

\begin{prop}\ollabel{prop:provability-lif}
  \begin{enumerate}
  \item \ollabel{prop:provability-lif-left}  $!A, !A \lif !B \Proves !B$.
  \item \ollabel{prop:provability-lif-right}
    Both $\lnot !A \Proves !A \lif !B$ and $!B \Proves !A \lif !B$.
  \end{enumerate}
\end{prop}

\begin{proof}
  \begin{enumerate}
  \item We can !!{derive}:
    \begin{prooftree}
      \AxiomC{$!A \lif !B$}
      \AxiomC{$!A$}
      \RightLabel{\Elim{\lif}}
      \BinaryInfC{$!B$}
    \end{prooftree}
    
  \item This is shown by the following two !!{derivation}s:
    \begin{prooftree}
      \AxiomC{$\lnot !A$}
      \AxiomC{$\Discharge{!A}{1}$}
      \RightLabel{\Elim{\lnot}}
      \BinaryInfC{$\lfalse$}
      \RightLabel{\FalseInt}
      \UnaryInfC{$!B$}
      \DischargeRule{\Intro{\lif}}{1}
      \UnaryInfC{$!A \lif !B$}
      \DisplayProof\qquad\bottomAlignProof
      \AxiomC{$!B$}
      \RightLabel{\Intro{\lif}}
      \UnaryInfC{$!A \lif !B$}
    \end{prooftree}
    Note that $\Intro{\lif}$ may, but does not have to, !!{discharge} the
    assumption~$!A$.
  \end{enumerate}
\end{proof}




  

\olfileid{fol}{ntd}{qpr}

\olsection{\usetoken{S}{derivability} and the Quantifiers}

\begin{explain}
  The completeness theorem also requires that the natural deduction
  rules yield the facts about~$\Proves$ established in this section.
\end{explain}

\begin{thm}
\ollabel{thm:strong-generalization} If $c$ is a constant not occurring
in $\Gamma$ or $!A(x)$ and $\Gamma \Proves !A(c)$, then $\Gamma
\Proves \lforall[x][!A(x)]$.
\end{thm}

\begin{proof}
Let $\delta$ be !!a{derivation} of $!A(c)$ from $\Gamma$.  By adding a
\Intro{\lforall} inference, we obtain !!a{derivation} of
$\lforall[x][!A(x)]$. Since $c$ does not occur in $\Gamma$ or $!A(x)$,
the eigenvariable condition is satisfied.
\end{proof}

\begin{prop}
\ollabel{prop:provability-quantifiers}
\begin{tagenumerate}{prvEx,prvAll}
$!A(t) \Proves \lexists[x][!A(x)]$.

$\lforall[x][!A(x)] \Proves !A(t)$.
\end{tagenumerate}
\end{prop}

\begin{proof}
\begin{tagenumerate}{prvEx,prvAll}
The following is !!a{derivation}
  of~$\lexists[x][!A(x)]$ from~$!A(t)$:
\begin{prooftree}
\AxiomC{$!A(t)$}
\RightLabel{\Intro{\lexists}}
\UnaryInfC{$\lexists[x][!A(x)]$}
\end{prooftree}

The following is !!a{derivation} of~$!A(t)$
  from~$\lforall[x][!A(x)]$:
\begin{prooftree}
\AxiomC{$\lforall[x][!A(x)]$}
\RightLabel{\Elim{\lforall}}
\UnaryInfC{$!A(t)$}
\end{prooftree}
\end{tagenumerate}
\end{proof}






\olfileid{fol}{ntd}{sou}
      
\olsection{Soundness}

\begin{explain}
!!^a{derivation} system, such as natural deduction, is \emph{sound}
if it cannot !!{derive} things that do not actually follow.  Soundness is
thus a kind of guaranteed safety property for !!{derivation} systems.
Depending on which proof theoretic property is in question, we would
like to know for instance, that
\begin{enumerate}
\item every !!{derivable} !!{sentence} is valid;
\item if !!a{sentence} is !!{derivable} from some others, it is also a
  consequence of them;
\item if a set of !!{sentence}s is inconsistent, it is unsatisfiable.
\end{enumerate}
These are important properties of !!a{derivation} system. If any of
them do not hold, the !!{derivation} system is deficient---it would
!!{derive} too much.  Consequently, establishing the soundness of a
!!{derivation} system is of the utmost importance.
\end{explain}

\begin{thm}[Soundness]
\ollabel{thm:soundness}
If $!A$ is !!{derivable} from the !!{undischarged} assumptions
$\Gamma$, then $\Gamma \Entails !A$.
\end{thm}

\begin{proof}
Let $\delta$ be !!a{derivation} of $!A$. We proceed by
induction on the number of inferences in~$\delta$.

For the induction basis we show the claim if the number of inferences
is~$0$. In this case, $\delta$ consists only of a single
!!{sentence}~$!A$, i.e., an assumption. That assumption is
!!{undischarged}, since assumptions can only be !!{discharged} by
inferences, and there are no inferences.  So, any
!!{structure}~$\Struct{M}$
that satisfies all of the !!{undischarged} assumptions of the proof
also satisfies~$!A$.

Now for the inductive step. Suppose that $\delta$ contains~$n$
inferences. The premise(s) of the lowermost inference are !!{derive}d
using sub-!!{derivation}s, each of which contains fewer than~$n$
inferences.  We assume the induction hypothesis: The premises of the
lowermost inference follow from the !!{undischarged} assumptions of
the sub-!!{derivation}s ending in those premises.  We have to show
that the conclusion~$!A$ follows from the !!{undischarged} assumptions
of the entire proof.

We distinguish cases according to the type of the lowermost inference.
First, we consider the possible inferences with only one premise.

\begin{enumerate}
\item Suppose that the last inference is \Intro{\lnot}: The
  !!{derivation} has the form
  \begin{prooftree}
    \AxiomC{$\Gamma, \Discharge{!A}{n}$}
    \RightLabel{$\delta_1$}
    \DeduceC{$\lfalse$}
    \DischargeRule{\Intro{\lnot}}{n}
    \UnaryInfC{$\lnot !A$}
  \end{prooftree}
  By inductive hypothesis, $\lfalse$ follows from the !!{undischarged}
  assumptions $\Gamma \cup \{!A\}$ of~$\delta_1$. Consider
  !!a{structure}~$\Struct{M}$. We
  need to show that, if $\Sat{M}{\Gamma}$, then
  $\Sat{M}{\lnot !A}$. Suppose for reductio
  that $\Sat{M}{\Gamma}$, but
  $\Sat/{M}{\lnot !A}$, i.e.,
  $\Sat{M}{!A}$. This would mean that
  $\Sat{M}{\Gamma \cup \{!A\}}$. This is
  contrary to our inductive hypothesis. So,
  $\Sat{M}{\lnot !A}$.

\item The last inference is \Elim{\land}: There are two variants: $!A$
  or $!B$ may be inferred from the premise $!A \land !B$. Consider the
  first case. The !!{derivation}~$\delta$ looks like this:
  \begin{prooftree}
    \AxiomC{$\Gamma$}
    \RightLabel{$\delta_1$}
    \DeduceC{$!A \land !B$}
    \RightLabel{\Elim{\land}}
    \UnaryInfC{$!A$}
  \end{prooftree}
  By inductive hypothesis, $!A \land !B$ follows from the
  !!{undischarged} assumptions~$\Gamma$ of~$\delta_1$. Consider
  !!a{structure}~$\Struct{M}$. We need to
  show that, if $\Sat{M}{\Gamma}$, then
  $\Sat{M}{!A}$. Suppose
  $\Sat{M}{\Gamma}$. By our inductive
  hypothesis ($\Gamma \Entails !A \land !B$), we know that
  $\Sat{M}{!A \land !B}$. By definition,
  $\Sat{M}{!A \land !B}$ iff
  $\Sat{M}{!A}$ and
  $\Sat{M}{!B}$.  (The case where $!B$ is
  inferred from $!A \land !B$ is handled similarly.)
  
\item The last inference is \Intro{\lor}: There are two variants: $!A
  \lor !B$ may be inferred from the premise~$!A$ or the
  premise~$!B$. Consider the first case. The !!{derivation} has the form
  \begin{prooftree}
    \AxiomC{$\Gamma$}
    \RightLabel{$\delta_1$}
    \DeduceC{$!A$}
    \RightLabel{\Intro{\lor}}
    \UnaryInfC{$!A \lor !B$}
  \end{prooftree}
  By inductive hypothesis, $!A$ follows from the !!{undischarged}
  assumptions~$\Gamma$ of~$\delta_1$.  Consider
  !!a{structure}~$\Struct{M}$. We
  need to show that, if $\Sat{M}{\Gamma}$, then
  $\Sat{M}{!A \lor !B}$. Suppose
  $\Sat{M}{\Gamma}$; then
  $\Sat{M}{!A}$ since $\Gamma \Entails !A$ (the
  inductive hypothesis). So it must also be the case that
  $\Sat{M}{!A \lor !B}$.  (The case where $!A
  \lor !B$ is inferred from~$!B$ is handled similarly.)
  
\item The last inference is \Intro{\lif}: $!A \lif !B$ is inferred
  from a subproof with assumption~$!A$ and conclusion~$!B$, i.e.,
  \begin{prooftree}
    \AxiomC{$\Gamma, \Discharge{!A}{n}$}
    \RightLabel{$\delta_1$}
    \DeduceC{$!B$}
    \DischargeRule{\Intro{\lif}}{n}
    \UnaryInfC{$!A \lif !B$}
  \end{prooftree}
  By inductive hypothesis, $!B$ follows from the !!{undischarged}
  assumptions of~$\delta_1$, i.e., $\Gamma \cup \{!A\} \Entails
  !B$. Consider
  !!a{structure}~$\Struct{M}$. The
  !!{undischarged} assumptions of~$\delta$ are just $\Gamma$, since
  $!A$ is discharged at the last inference. So we need to show that
  $\Gamma \Entails !A \lif !B$.  For reductio, suppose that for some
  !!{structure}~$\Struct{M}$,
  $\Sat{M}{\Gamma}$ but
  $\Sat/{M}{!A \lif !B}$. So,
  $\Sat{M}{!A}$ and
  $\Sat/{M}{!B}$. But by hypothesis, $!B$ is a
  consequence of $\Gamma \cup \{!A\}$, i.e.,
  $\Sat{M}{!B}$, which is a contradiction. So,
  $\Gamma \Entails !A \lif !B$.

\item The last inference is \FalseInt: Here, $\delta$ ends in
  \begin{prooftree}
    \AxiomC{$\Gamma$}
    \RightLabel{$\delta_1$}
    \DeduceC{$\lfalse$}
    \RightLabel{\FalseInt}
    \UnaryInfC{$!A$}
  \end{prooftree}
  By induction hypothesis, $\Gamma \Entails \lfalse$. We have to show
  that $\Gamma \Entails !A$. Suppose not; then for
  some~$\Struct{M}$ we have
    $\Sat{M}{\Gamma}$ and
    $\Sat/{M}{!A}$. But we always have
    $\Sat/{M}{\lfalse}$, so this would mean
    that $\Gamma \Entails/ \lfalse$, contrary to the induction
    hypothesis.

\item The last inference is \FalseCl: Exercise.


\item The last inference is \Intro{\lforall}: Then $\delta$ has the form
  \begin{prooftree}
    \AxiomC{$\Gamma$}
    \RightLabel{$\delta_1$}
    \DeduceC{$!A(a)$}
    \RightLabel{\Intro{\lforall}}
    \UnaryInfC{$\lforall[x][!A(x)]$}
  \end{prooftree}
  The premise $!A(a)$ is a consequence of the !!{undischarged}
  assumptions $\Gamma$ by induction hypothesis.  Consider some
  structure, $\Struct{M}$, such that $\Sat{M}{\Gamma}$.  We need to
  show that $\Sat{M}{\lforall[x][!A(x)]}$. Since $\lforall[x][!A(x)]$
  is !!a{sentence}, this means we have to show that for every variable
  assignment~$s$, $\Sat{M}{!A(x)}[s]$
  (\olref[syn][ass]{prop:sat-quant}). Since $\Gamma$ consists entirely
  of sentences, $\Sat{M}{!B}[s]$ for all $!B \in \Gamma$ by
  \olref[syn][sat]{defn:satisfaction}.  Let $\Struct{M'}$ be like
  $\Struct{M}$ except that $\Assign{a}{M'} = s(x)$.  Since $a$ does
  not occur in~$\Gamma$, $\Sat{M'}{\Gamma}$ by
  \olref[syn][ext]{cor:extensionality-sent}. Since $\Gamma \Entails
  !A(a)$, $\Sat{M'}{!A(a)}$.  Since $!A(a)$ is !!a{sentence},
  $\Sat{M'}{!A(a)}[s]$ by
  \olref[syn][ass]{prop:sentence-sat-true}. $\Sat{M'}{!A(x)}[s]$ iff
  $\Sat{M'}{!A(a)}$ by \olref[syn][ext]{prop:ext-formulas} (recall
  that $!A(a)$ is just $\Subst{!A(x)}{a}{x}$). So,
  $\Sat{M'}{!A(x)}[s]$. Since $a$ does not occur in~$!A(x)$, by
  \olref[syn][ext]{prop:extensionality}, $\Sat{M}{!A(x)}[s]$. But $s$
  was an arbitrary variable assignment, so
  $\Sat{M}{\lforall[x][!A(x)]}$.
  
\item The last inference is \Intro{\lexists}: Exercise.

\item The last inference is \Elim{\forall}: Exercise.

\end{enumerate}

Now let's consider the possible inferences with several premises:
\Elim{\lor}, \Intro{\land}, \Elim{\lif}, and
  \Elim{\lexists}.
\begin{enumerate}

\item The last inference is \Intro{\land}. $!A \land !B$ is inferred
  from the premises $!A$ and $!B$ and $\delta$ has the form
  \begin{prooftree}
    \AxiomC{$\Gamma_1$}
    \RightLabel{$\delta_1$}
    \DeduceC{$!A$}
    \AxiomC{$\Gamma_2$}
    \RightLabel{$\delta_2$}
    \DeduceC{$!B$}
    \RightLabel{\Intro{\land}}
    \BinaryInfC{$!A \land !B$}
  \end{prooftree}
  By induction hypothesis, $!A$ follows from the !!{undischarged}
  assumptions~$\Gamma_1$ of~$\delta_1$ and $!B$ follows from the
  !!{undischarged} assumptions~$\Gamma_2$ of~$\delta_2$. The
  !!{undischarged} assumptions of~$\delta$ are $\Gamma_1 \cup
  \Gamma_2$, so we have to show that $\Gamma_1 \cup \Gamma_2 \Entails
  !A \land !B$. Consider
  !!a{structure}~$\Struct{M}$
  with $\Sat{M}{\Gamma_1 \cup \Gamma_2}$. Since
  $\Sat{M}{\Gamma_1}$, it must be the case that
  $\Sat{M}{!A}$ as $\Gamma_1 \Entails !A$, and
  since $\Sat{M}{\Gamma_2}$,
  $\Sat{M}{!B}$ since $\Gamma_2 \Entails
  !B$. Together, $\Sat{M}{!A \land !B}$.
  
\item The last inference is \Elim{\lor}: Exercise.

\item The last inference is \Elim{\lif}. $!B$ is inferred from the
  premises $!A \lif !B$ and~$!A$. The !!{derivation}~$\delta$ looks like this:
  \begin{prooftree}
    \AxiomC{$\Gamma_1$}
    \RightLabel{$\delta_1$}
    \DeduceC{$!A \lif !B$}
    \AxiomC{$\Gamma_2$}
    \RightLabel{$\delta_2$}
    \DeduceC{$!A$}
    \RightLabel{\Elim{\lif}}
    \BinaryInfC{$!B$}
  \end{prooftree}
  By induction hypothesis, $!A \lif !B$ follows from the
  !!{undischarged} assumptions~$\Gamma_1$ of~$\delta_1$ and $!A$
  follows from the !!{undischarged} assumptions~$\Gamma_2$
  of~$\delta_2$. Consider
  !!a{structure}~$\Struct{M}$. We
  need to show that, if $\Sat{M}{\Gamma_1 \cup
    \Gamma_2}$, then $\Sat{M}{!B}$. Suppose
  $\Sat{M}{\Gamma_1 \cup \Gamma_2}$. Since
  $\Gamma_1 \Entails !A \lif !B$, $\Sat{M}{!A
    \lif !B}$. Since $\Gamma_2 \Entails !A$, we have
  $\Sat{M}{!A}$. This means that
  $\Sat{M}{!B}$ (For if
  $\Sat/{M}{!B}$, since
  $\Sat{M}{!A}$, we'd have
  $\Sat/{M}{!A \lif !B}$,
  contradicting~$\Sat{M}{!A \lif !B}$).

\item The last inference is \Elim{\lnot}: Exercise.
  
The last inference is \Elim{\lexists}: Exercise.
\end{enumerate}
\end{proof}


\begin{prob}
Complete the proof of \olref[fol][ntd][sou]{thm:soundness}.
\end{prob}




\begin{cor}
\ollabel{cor:weak-soundness}
If $\Proves !A$, then $!A$ is valid.
\end{cor}

\begin{cor}
\ollabel{cor:consistency-soundness}
If $\Gamma$ is satisfiable, then it is consistent.
\end{cor}

\begin{proof}
We prove the contrapositive.  Suppose that $\Gamma$ is not consistent.
Then $\Gamma \Proves \lfalse$, i.e., there is !!a{derivation} of
$\lfalse$ from !!{undischarged} assumptions in~$\Gamma$. By
\olref{thm:soundness}, any
!!{structure}~$\Struct{M}$
that satisfies $\Gamma$ must satisfy $\lfalse$.  Since
$\Sat/{M}{\lfalse}$ for every
!!{structure}~$\Struct{M}$,
no $\Struct{M}$ can satisfy $\Gamma$,
i.e., $\Gamma$ is not satisfiable.
\end{proof}




  

\olfileid{fol}{ntd}{ide}

\olsection{\usetoken{P}{derivation} with \usetoken{S}{identity}}

!!^{derivation}s with !!{identity} require additional inference rules.

\begin{defish}
\AxiomC{}
\RightLabel{\Intro{\eq}}
\UnaryInfC{$\eq[t][t]$}
\DisplayProof
\hfill
\begin{tabular}{r}
\AxiomC{$\eq[t_1][t_2]$}
\AxiomC{$!A(t_1)$}
\RightLabel{\Elim{\eq}}
\BinaryInfC{$!A(t_2)$}
\DisplayProof
\\[3ex]
\AxiomC{$\eq[t_1][t_2]$}
\AxiomC{$!A(t_2)$}
\RightLabel{\Elim{\eq}}
\BinaryInfC{$!A(t_1)$}
\DisplayProof
\end{tabular}
\end{defish}

In the above rules, $t$, $t_1$, and $t_2$ are closed terms. The
\Intro{\eq} rule allows us to !!{derive} any identity statement of the
form $\eq[t][t]$ outright, from no assumptions.

\begin{ex}
If $s$ and $t$ are closed terms, then $!A(s), \eq[s][t] \Proves !A(t)$:
\begin{prooftree}
\AxiomC{$\eq[s][t]$}
\AxiomC{$!A(s)$}
\RightLabel{$\Elim{\eq}$}
\BinaryInfC{$!A(t)$}
\end{prooftree}
This may be familiar as the ``principle of substitutability of
identicals,'' or Leibniz' Law.
\end{ex}

\begin{prob}
Prove that $=$ is both symmetric and transitive, i.e., give
!!{derivation}s of $\lforall[x][\lforall[y][(\eq[x][y] \lif
    \eq[y][x])]]$ and $\lforall[x][\lforall[y][\lforall[z]((\eq[x][y]
    \land \eq[y][z]) \lif \eq[x][z])]]$
\end{prob}

\begin{ex}
We !!{derive} the !!{sentence}
\begin{align*}
& \lforall[x][\lforall[y][((!A(x) \land !A(y)) \lif \eq[x][y])]]
\intertext{from the !!{sentence}}
& \lexists[x][\lforall[y][(!A(y) \lif \eq[y][x])]]
\end{align*}
We develop the !!{derivation} backwards:
\begin{prooftree}
\AxiomC{$\lexists[x][\lforall[y][(!A(y) \lif \eq[y][x])]]
\quad \Discharge{!A(a) \land !A(b)}{1}$}
\DeduceC{$\eq[a][b]$}
\DischargeRule{\Intro{\lif}}{1}
\UnaryInfC{$((!A(a) \land !A(b)) \lif \eq[a][b])$}
\RightLabel{\Intro{\lforall}}
\UnaryInfC{$\lforall[y][((!A(a) \land !A(y)) \lif \eq[a][y])]$}
\RightLabel{\Intro{\lforall}}
\UnaryInfC{$\lforall[x][\lforall[y][((!A(x) \land !A(y)) \lif \eq[x][y])]]$}
\end{prooftree}
We'll now have to use the main assumption: since it is an existential
!!{formula}, we use \Elim{\lexists} to !!{derive} the intermediary
conclusion $\eq[a][b]$.
\begin{prooftree}
\AxiomC{$\lexists[x][\lforall[y][(!A(y) \lif \eq[y][x])]]$}
\AxiomC{$\Discharge{\lforall[y][(!A(y) \lif \eq[y][c]])}{2}$}
\noLine
\UnaryInfC{$\Discharge{!A(a) \land !A(b)}{1}$}
\DeduceC{$\eq[a][b]$}
\DischargeRule{\Elim{\lexists}}{2}
\BinaryInfC{$\eq[a][b]$}
\DischargeRule{\Intro{\lif}}{1}
\UnaryInfC{$((!A(a) \land !A(b)) \lif \eq[a][b])$}
\RightLabel{\Intro{\lforall}}
\UnaryInfC{$\lforall[y][((!A(a) \land !A(y)) \lif \eq[a][y])]$}
\RightLabel{\Intro{\lforall}}
\UnaryInfC{$\lforall[x][\lforall[y][((!A(x) \land !A(y)) \lif \eq[x][y])]]$}
\end{prooftree}
The sub-!!{derivation} on the top right is completed by using its
assumptions to show that $\eq[a][c]$ and $\eq[b][c]$. This requires two
separate !!{derivation}s. The !!{derivation} for $\eq[a][c]$ is as
follows:
\begin{prooftree}
\AxiomC{$\Discharge{\lforall[y][(!A(y) \lif \eq[y][c]])}{2}$}
\RightLabel{\Elim{\lforall}}
\UnaryInfC{$!A(a) \lif \eq[a][c]$}
\AxiomC{$\Discharge{!A(a) \land !A(b)}{1}$}
\RightLabel{\Elim{\land}}
\UnaryInfC{$!A(a)$}
\RightLabel{\Elim{\lif}}
\BinaryInfC{$\eq[a][c]$}
\end{prooftree}
From $\eq[a][c]$ and $\eq[b][c]$ we !!{derive} $\eq[a][b]$ by
\Elim{\eq}.
\end{ex}

\begin{prob}
Give !!{derivation}s of the following !!{formula}s:
\begin{enumerate}
\item $\lforall[x][\lforall[y][((\eq[x][y] \land !A(x)) \lif !A(y))]]$
\item $\lexists[x][!A(x)] \land \lforall[y][\lforall[z][((!A(y) \land
    !A(z)) \lif \eq[y][z])]] \lif \lexists[x][(!A(x) \land
  \lforall[y][(!A(y) \lif \eq[y][x])])]$
\end{enumerate}
\end{prob}


  

\olfileid{fol}{ntd}{sid}

\olsection{Soundness with \usetoken{S}{identity}}

\begin{prop}
Natural deduction with rules for $\eq$ is sound.
\end{prop}

\begin{proof}
Any !!{formula} of the form $\eq[t][t]$ is valid, since
for every !!{structure}~$\Struct M$, $\Sat{M}{\eq[t][t]}$. (Note that
we assume the term $t$ to be closed, i.e., it contains no variables,
so variable assignments are irrelevant).

Suppose the last inference in !!a{derivation} is \Elim{\eq}, i.e., the
!!{derivation} has the following form:
\begin{prooftree}
  \AxiomC{$\Gamma_1$}
  \RightLabel{$\delta_1$}
  \DeduceC{$\eq[t_1][t_2]$}
  \AxiomC{$\Gamma_2$}
  \RightLabel{$\delta_2$}
  \DeduceC{$!A(t_1)$}
  \RightLabel{\Elim{\eq}}
  \BinaryInfC{$!A(t_2)$}
\end{prooftree}
The premises $\eq[t_1][t_2]$ and $!A(t_1)$ are !!{derive}d from
!!{undischarged} assumptions~$\Gamma_1$ and $\Gamma_2$, respectively.
We want to show that $!A(t_2)$ follows from $\Gamma_1 \cup \Gamma_2$.
Consider !!a{structure}~$\Struct{M}$ with $\Sat{M}{\Gamma_1 \cup
  \Gamma_2}$. By induction hypothesis, $\Sat{M}{!A(t_1)}$ and
$\Sat{M}{\eq[t_1][t_2]}$. Therefore, $\Value{t_1}{M} = \Value{t_2}{M}$. Let
$s$ be any variable assignment, and $m = \Value{t_1}{M} = \Value{t_2}{M}$. By
\olref[fol][syn][ext]{prop:ext-formulas}, $\Sat{M}{!A(t_1)}[s]$ iff
$\Sat{M}{!A(x)}[\Subst{s}{m}{x}]$ iff $\Sat{M}{!A(t_2)}[s]$. Since
$\Sat{M}{!A(t_1)}$, we have $\Sat{M}{!A(t_2)}$.
\end{proof}




\OLEndChapterHook





  

